\section{Conclusion}
\label{sec:conclusion}

In this paper, we have shown that for every diagram in which each $c$-component has at most two hyperedges, and across every causal order and positive table, the reduced partial identification program for a linear query reduces to one linear program returning reproducibility and the sharp bounds if and only if the three tests of \cref{thm:lp} hold. Linear assumptions, as in \cref{prop:assumptions}, preserve that conclusion: the reduced program stays a linear program, and the optimum remains the identified set under the assumptions. We introduce a new software package, \pid{}, that reads the tests off the diagram and the query, returns the interval, and records whether the interval is the identified set.

We have also shown that a component with three or more hyperedges lies outside the equivalence. We do not treat continuous variables, symbolic closed forms, sharp product independence, unanchored Rosenbaum $\Gamma$, raw mutual-information budgets, and free-denominator queries at leftover degree two or more. What remains open is a sharp solver for fractional objectives when two simplices still multiply, for independence splits, and for nested optima on every bilinear instance, together with the graphical question for $c$-components with three or more hyperedges. 
